\documentclass[a4paper,12pt]{article}

% Pakete
\usepackage[utf8]{inputenc}
\usepackage[T1]{fontenc}
\usepackage[ngerman]{babel}
\usepackage{geometry}
\usepackage{graphicx}
\usepackage{xcolor}
\usepackage{listings}
\usepackage{enumitem}
\usepackage{booktabs}
\usepackage{tabularx}
\usepackage{hyperref}
\usepackage{fancyhdr}
\usepackage{titlesec}
\usepackage{tcolorbox}
\usepackage{multicol}

\geometry{left=2.5cm, right=2.5cm, top=2.5cm, bottom=2.5cm}

% Farben
\definecolor{primary}{HTML}{1976D2}
\definecolor{success}{HTML}{4CAF50}
\definecolor{warning}{HTML}{FF9800}
\definecolor{danger}{HTML}{F44336}
\definecolor{dark}{HTML}{1E1E1E}
\definecolor{codebg}{HTML}{F5F5F5}

% Hyperlinks
\hypersetup{
    colorlinks=true,
    linkcolor=primary,
    urlcolor=primary
}

% Header/Footer
\pagestyle{fancy}
\fancyhf{}
\fancyhead[L]{\textcolor{primary}{\textbf{MaterialManager R03}}}
\fancyhead[R]{\textcolor{gray}{Technische Dokumentation}}
\fancyfoot[C]{\thepage}
\renewcommand{\headrulewidth}{1pt}
\renewcommand{\headrule}{\hbox to\headwidth{\color{primary}\leaders\hrule height \headrulewidth\hfill}}

% Code-Stil
\lstset{
    backgroundcolor=\color{codebg},
    basicstyle=\ttfamily\small,
    breaklines=true,
    frame=single,
    rulecolor=\color{gray!30},
    numbers=left,
    numberstyle=\tiny\color{gray},
    keywordstyle=\color{primary}\bfseries,
    stringstyle=\color{success},
    commentstyle=\color{gray}\itshape
}

% Info-Box
\newtcolorbox{infobox}[1][]{
    colback=primary!5,
    colframe=primary,
    fonttitle=\bfseries,
    title=#1
}

\newtcolorbox{warnbox}[1][]{
    colback=warning!5,
    colframe=warning,
    fonttitle=\bfseries,
    title=#1
}

\newtcolorbox{successbox}[1][]{
    colback=success!5,
    colframe=success,
    fonttitle=\bfseries,
    title=#1
}

\begin{document}

% ==========================================
% TITELSEITE
% ==========================================
\begin{titlepage}
\centering
\vspace*{3cm}

{\Huge\bfseries\textcolor{primary}{MaterialManager R03}}\\[0.5cm]
{\Large\textcolor{gray}{Technische Dokumentation \& Bewertung}}\\[2cm]

\begin{tcolorbox}[colback=dark,colframe=dark,coltext=white,width=12cm,halign=center]
\centering
{\large\bfseries Materialverwaltung für Blechlager}\\[0.3cm]
{\normalsize WPF Desktop-Anwendung $\bullet$ .NET 8 $\bullet$ C\# 12}\\[0.3cm]
{\small Dark Theme $\bullet$ Auto-Save $\bullet$ ERP-Schnittstelle $\bullet$ Netzwerk-Modus}
\end{tcolorbox}

\vspace{2cm}

\begin{tabular}{ll}
\textbf{Version:} & 1.0.0 \\
\textbf{Plattform:} & Windows 10/11 (x64) \\
\textbf{Framework:} & .NET 8.0 (Self-Contained) \\
\textbf{Lizenz:} & Demo (30 Tage) \\
\textbf{Stand:} & \today \\
\end{tabular}

\vfill
{\small Copyright \copyright\ 2025 MaterialManager}
\end{titlepage}

% ==========================================
% INHALTSVERZEICHNIS
% ==========================================
\tableofcontents
\newpage

% ==========================================
% 1. PROGRAMMÜBERSICHT
% ==========================================
\section{Programmübersicht}

\subsection{Zweck}
MaterialManager R03 ist eine professionelle \textbf{Materialverwaltungssoftware} für industrielle Blechlager. Die Software verwaltet den gesamten Lebenszyklus von Blechtafeln — vom Wareneingang über die Lagerung bis zum Verbrauch und der Restverwertung.

\subsection{Zielgruppe}
\begin{itemize}
    \item Metallbau-Unternehmen
    \item Blechbearbeitungsbetriebe
    \item Stahlhandel
    \item Schlossereien und Schweißbetriebe
\end{itemize}

\subsection{Technologie-Stack}

\begin{table}[h]
\centering
\begin{tabularx}{\textwidth}{l X}
\toprule
\textbf{Komponente} & \textbf{Technologie} \\
\midrule
Sprache & C\# 12.0 \\
Framework & .NET 8.0 (LTS) \\
UI-Framework & WPF (Windows Presentation Foundation) \\
Datenspeicherung & Excel (.xlsx) via ClosedXML \\
Deployment & Self-Contained (keine Installation nötig) \\
Architektur & MVVM-ähnlich (Code-Behind) \\
\bottomrule
\end{tabularx}
\caption{Technologie-Stack}
\end{table}

% ==========================================
% 2. BEWERTUNG
% ==========================================
\section{Bewertung des Programms}

\subsection{Gesamtbewertung}

\begin{successbox}[Gesamtnote: 8.5 / 10 — Sehr gut für eine Branchen-Lösung]
Die Software deckt die wesentlichen Anforderungen einer industriellen Materialverwaltung ab und bietet darüber hinaus Funktionen, die in vielen kommerziellen Produkten fehlen (Verschnitt-Optimierung, automatische Regalzuweisung).
\end{successbox}

\subsection{Bewertung nach Kategorien}

\begin{table}[h]
\centering
\begin{tabularx}{\textwidth}{l c X}
\toprule
\textbf{Kategorie} & \textbf{Note} & \textbf{Begründung} \\
\midrule
Funktionsumfang & 9/10 & 20+ Features, deckt gesamten Workflow ab \\
Benutzeroberfläche & 8/10 & Modernes Dark Theme, intuitive Bedienung \\
Datensicherheit & 8/10 & Auto-Save, Backup, Undo-Funktion \\
ERP-Integration & 8/10 & SAP, SAGE, BMEcat, Universal CSV \\
Netzwerkfähigkeit & 7/10 & File-basiert mit Locking (kein SQL-Server) \\
Performance & 9/10 & Schnell, kein spürbarer Lag \\
Code-Qualität & 7/10 & Funktional gut, könnte mehr MVVM nutzen \\
Wartbarkeit & 7/10 & Service-Architektur, aber Code-Behind lastig \\
Dokumentation & 6/10 & README vorhanden, API-Doku fehlt \\
\bottomrule
\end{tabularx}
\caption{Bewertung nach Kategorien}
\end{table}

\subsection{Stärken}

\begin{enumerate}
    \item \textbf{Vollständiger Workflow:} Material anlegen $\rightarrow$ Lagern $\rightarrow$ Verbrauchen $\rightarrow$ Reste verwalten
    \item \textbf{Automatische Regalzuweisung:} Intelligente Zuordnung basierend auf Material, Legierung, Form und Stärke
    \item \textbf{Verschnitt-Optimierung:} Findet das beste Material für ein benötigtes Teil (Reste bevorzugt)
    \item \textbf{Auto-Save:} Kein manuelles Speichern nötig, keine Datenverluste
    \item \textbf{ERP-Kompatibilität:} SAP MM, SAGE, BMEcat — sofort nutzbar
    \item \textbf{Self-Contained:} Keine .NET-Installation nötig, USB-Stick-fähig
    \item \textbf{Netzwerk-Modus:} Mehrere PCs können gleichzeitig arbeiten
\end{enumerate}

\subsection{Verbesserungspotenzial}

\begin{enumerate}
    \item \textbf{Datenbank:} Excel als Backend hat Grenzen bei großen Datenmengen ($>$10.000 Positionen)
    \item \textbf{Benutzer-Verwaltung:} Kein Login-System, keine Rollen (Admin/Lager/Büro)
    \item \textbf{Audit-Trail:} Keine vollständige Nachverfolgung aller Änderungen (ISO 9001)
    \item \textbf{MVVM:} Code-Behind könnte durch ViewModels ersetzt werden
    \item \textbf{Unit-Tests:} Keine automatisierten Tests vorhanden
\end{enumerate}

% ==========================================
% 3. FUNKTIONSÜBERSICHT
% ==========================================
\section{Funktionsübersicht}

\subsection{Kernfunktionen}

\begin{table}[h]
\centering
\begin{tabularx}{\textwidth}{l l X}
\toprule
\textbf{Funktion} & \textbf{Button} & \textbf{Beschreibung} \\
\midrule
Material anlegen & Material neu & Stahl/Edelstahl/Aluminium mit allen Eigenschaften \\
Material bearbeiten & Doppelklick & Alle Felder editierbar, Änderungsdatum automatisch \\
Material löschen & Ausgewählte löschen & Mehrfachauswahl per Checkbox möglich \\
Tafel verbrauchen & Tafel verbrauchen & Stückzahl $-1$, Reste mit Maßen anlegen \\
Verschnitt-Check & Verschnitt-Check & Findet bestes Material für benötigte Maße \\
Reste suchen & Reste suchen & Suche nach passendem Rest-Material \\
Rückgängig & $\hookleftarrow$ Rückgängig & Letzte Löschung rückgängig machen \\
\bottomrule
\end{tabularx}
\caption{Kernfunktionen}
\end{table}

\subsection{Verwaltungsfunktionen}

\begin{table}[h]
\centering
\begin{tabularx}{\textwidth}{l X}
\toprule
\textbf{Funktion} & \textbf{Beschreibung} \\
\midrule
Inventur & Soll/Ist-Vergleich aller Regale \\
Lagerkarte & Visuelle Regalübersicht mit Ampelfarben \\
Statistik & Auswertungen und Übersichten \\
Drucken & Lagerliste, Resteliste, Etiketten mit QR-Code \\
ERP-Export & SAP MM, SAGE, BMEcat XML, Universal CSV \\
Netzwerk & Gemeinsamer Datenzugriff über Netzwerkpfad \\
Excel Export/Import & Datenaustausch via .xlsx \\
\bottomrule
\end{tabularx}
\caption{Verwaltungsfunktionen}
\end{table}

\subsection{Dashboard-Widgets}

Das Hauptfenster zeigt vier Karten:

\begin{enumerate}
    \item \textbf{Gesamtgewicht:} Summe aller Materialien in kg
    \item \textbf{Regalauslastung:} Durchschnittliche Auslastung aller Regale in \%
    \item \textbf{EU Palette:} Auslastung der Rest-Palette (klickbar für Details)
    \item \textbf{Niedrige Bestände:} Warnungen für Materialien unter 3 Tafeln
\end{enumerate}

% ==========================================
% 4. ARCHITEKTUR
% ==========================================
\section{Software-Architektur}

\subsection{Projektstruktur}

\begin{lstlisting}[language={},title=Projektstruktur]
MaterialManager_R03/
|-- Models/
|   |-- MaterialItem.cs          (Datenmodell)
|   |-- MaterialDefinitions.cs   (Stammdaten)
|   |-- GroupedMaterialItem.cs    (Gruppierung)
|-- Views/
|   |-- MaterialDialog.xaml(.cs)  (Material anlegen/bearbeiten)
|   |-- TafelVerbrauchDialog.cs   (Tafel verbrauchen)
|   |-- ResteSucheDialog.xaml(.cs)(Reste suchen)
|   |-- StatistikDialog.cs       (Statistik)
|   |-- InventurDialog.cs        (Inventur)
|   |-- BenachrichtigungsDialog.cs
|   |-- NiedrigeBestaendeDialog.xaml(.cs)
|   |-- EuPaletteDialog.xaml(.cs)
|-- Services/
|   |-- RegalService.cs          (Regal-Logik, 10 Regale)
|   |-- LagerService.cs          (Lager-Auslastung)
|   |-- ExcelService.cs          (Import/Export)
|   |-- VerschnittService.cs     (Verschnitt-Optimierung)
|   |-- ErpExportService.cs      (SAP/SAGE/BMEcat)
|   |-- NetzwerkService.cs       (Multi-PC Zugriff)
|   |-- BackupService.cs         (Auto-Backup alle 30 Min)
|   |-- BuchungsService.cs       (Ein-/Ausgangs-Protokoll)
|   |-- DruckService.cs          (Drucken & PDF)
|   |-- QrCodeService.cs         (Etiketten)
|   |-- UndoService.cs           (Ruckgaengig)
|   |-- BestellService.cs        (Bestellvorschlaege)
|   |-- LicenseService.cs        (30-Tage Demo)
|-- Converters/
|   |-- PercentToBrushConverter.cs
|   |-- StaerkeFormatConverter.cs
|-- MainWindow.xaml(.cs)         (Hauptfenster)
|-- App.xaml(.cs)                (App-Start)
\end{lstlisting}

\subsection{Regallogik}

Das System verwaltet \textbf{10 Regale + 1 EU-Palette} mit intelligenter Zuordnung:

\begin{table}[h]
\centering
\begin{tabularx}{\textwidth}{l r X}
\toprule
\textbf{Regal} & \textbf{Kapazität} & \textbf{Zuordnungsregel} \\
\midrule
Regal A & 800 kg & MF (Mittelformat) \\
Regal B & 1.500 kg & Aluminium/Edelstahl $\geq$ 6mm \\
Regal C & 1.000 kg & GF -- S355+ Stähle \\
Regal D & 1.000 kg & GF -- S235 und niedrigere Stähle \\
Regal E & 1.000 kg & GF -- Sonstige \\
Regal F & 800 kg & MF \\
Regal G--I & je 600 kg & KF (Kleinformat) \\
Regal J & 800 kg & MF \\
EU Palette & 2.000 kg & Kleine Materialien ($\leq$1500$\times$1300mm) \\
\bottomrule
\end{tabularx}
\caption{Regal-Zuordnung}
\end{table}

\subsection{Datenspeicherung}

\begin{infobox}[Speicher-Konzept]
\begin{itemize}
    \item \textbf{Lokal:} \texttt{\%LOCALAPPDATA\%\textbackslash MaterialManager\_R02\textbackslash autosave.xlsx}
    \item \textbf{Netzwerk:} \texttt{\textbackslash\textbackslash SERVER\textbackslash Freigabe\textbackslash materialbestand.xlsx}
    \item \textbf{Backup:} Automatisch alle 30 Minuten
    \item \textbf{Auto-Save:} Nach jeder Änderung (1 Sekunde Verzögerung)
\end{itemize}
\end{infobox}

% ==========================================
% 5. MATERIALTYPEN
% ==========================================
\section{Unterstützte Materialien}

\subsection{Materialarten}

\begin{table}[h]
\centering
\begin{tabularx}{\textwidth}{l l X}
\toprule
\textbf{Material} & \textbf{Dichte} & \textbf{Legierungen} \\
\midrule
Stahl & 7.850 kg/m³ & S235, S355, DC01, DC04, DD11, S460, HB400, HB500 \\
Edelstahl & 8.000 kg/m³ & 1.4301, 1.4404, 1.4571 \\
Aluminium & 2.700 kg/m³ & AlMg3, AlMg4.5Mn, AlMgSi1 \\
\bottomrule
\end{tabularx}
\caption{Materialarten und Legierungen}
\end{table}

\subsection{Formate}

\begin{table}[h]
\centering
\begin{tabular}{l l l}
\toprule
\textbf{Form} & \textbf{Standardmaß} & \textbf{Beschreibung} \\
\midrule
GF & 2500$\times$1250 mm & Großformat-Tafel \\
MF & 2000$\times$1000 mm & Mittelformat-Tafel \\
KF & 1000$\times$500 mm & Kleinformat-Tafel \\
Rest & frei & Reste aus Verbrauch \\
\bottomrule
\end{tabular}
\caption{Tafel-Formate}
\end{table}

% ==========================================
% 6. ERP-SCHNITTSTELLE
% ==========================================
\section{ERP-Schnittstelle}

\subsection{Unterstützte Formate}

\begin{table}[h]
\centering
\begin{tabularx}{\textwidth}{l l X}
\toprule
\textbf{Format} & \textbf{Dateitype} & \textbf{Einsatz} \\
\midrule
SAP MM & CSV (Semikolon) & SAP Materialwirtschaft (MATNR, MAKTX, WRKST) \\
SAGE & CSV (Semikolon) & SAGE 50/100/b7 \\
BMEcat & XML & Deutscher Industrie-Standard (Produktdatenaustausch) \\
Universal & CSV (Semikolon) & Jedes beliebige ERP-System \\
\bottomrule
\end{tabularx}
\caption{ERP-Export-Formate}
\end{table}

\subsection{SAP-Felder}

Der SAP-Export enthält folgende Felder:

\begin{multicols}{2}
\begin{itemize}[noitemsep]
    \item MATNR -- Materialnummer
    \item MAKTX -- Kurztext
    \item MATKL -- Materialklasse
    \item MEINS -- Mengeneinheit
    \item BRGEW -- Bruttogewicht
    \item GEWEI -- Gewichtseinheit
    \item LGORT -- Lagerort
    \item CHARG -- Charge/Restnummer
    \item LFDAT -- Lieferdatum
    \item LIFNR -- Lieferant
    \item EBELN -- Bestellnummer
    \item WRKST -- Werkstoffnummer (DIN)
\end{itemize}
\end{multicols}

\subsection{Werkstoffnummern-Zuordnung}

\begin{table}[h]
\centering
\begin{tabular}{l l l}
\toprule
\textbf{Legierung} & \textbf{Werkstoff-Nr.} & \textbf{Norm} \\
\midrule
S235 & 1.0037 & EN 10025 \\
S355 & 1.0577 & EN 10025 \\
DC01 & 1.0330 & EN 10130 \\
DC04 & 1.0338 & EN 10130 \\
1.4301 & 1.4301 & EN 10088 (V2A) \\
1.4404 & 1.4404 & EN 10088 (V4A) \\
1.4571 & 1.4571 & EN 10088 \\
\bottomrule
\end{tabular}
\caption{Werkstoffnummern nach DIN/EN}
\end{table}

% ==========================================
% 7. NETZWERK
% ==========================================
\section{Netzwerk-Betrieb}

\subsection{Funktionsweise}

\begin{infobox}[Multi-PC-Betrieb]
Alle PCs speichern und laden die gleiche Excel-Datei auf einem gemeinsamen Netzwerkpfad. Ein File-Locking-Mechanismus verhindert gleichzeitiges Überschreiben.
\end{infobox}

\subsection{Einrichtung}

\begin{enumerate}
    \item Netzwerkfreigabe erstellen: \texttt{\textbackslash\textbackslash SERVER\textbackslash MaterialManager}
    \item Auf jedem PC: Button \textbf{Netzwerk} $\rightarrow$ Pfad eintragen $\rightarrow$ \textbf{Aktivieren}
    \item Benutzername eintragen (zur Identifikation)
    \item Fertig — alle PCs arbeiten auf der gleichen Datei
\end{enumerate}

\subsection{Sicherheit}

\begin{itemize}
    \item \textbf{File-Lock:} Automatische Sperre bei Schreibzugriff
    \item \textbf{Auto-Release:} Lock wird nach 5 Minuten automatisch freigegeben
    \item \textbf{Fallback:} Bei Netzwerkausfall wird lokal gespeichert
    \item \textbf{Backup:} Automatisches Backup alle 30 Minuten
\end{itemize}

% ==========================================
% 8. STATISTIK
% ==========================================
\section{Code-Statistik}

\begin{table}[h]
\centering
\begin{tabular}{l r}
\toprule
\textbf{Metrik} & \textbf{Wert} \\
\midrule
Quelldateien (C\#) & 28 \\
XAML-Dateien & 7 \\
Services & 13 \\
Views / Dialoge & 9 \\
Models & 3 \\
Converter & 2 \\
Geschätzte LOC & $\sim$4.500 \\
NuGet-Pakete & 1 (ClosedXML) \\
\bottomrule
\end{tabular}
\caption{Code-Statistik}
\end{table}

% ==========================================
% 9. LIZENZ
% ==========================================
\section{Lizenzierung}

\subsection{Demo-Version}

\begin{warnbox}[Demo-Einschränkungen]
\begin{itemize}
    \item Laufzeit: \textbf{30 Tage} ab Erststart
    \item Alle Funktionen verfügbar (keine Einschränkungen)
    \item Countdown im Fenstertitel sichtbar
    \item Nach Ablauf: Programm startet nicht mehr
\end{itemize}
\end{warnbox}

\subsection{Technische Umsetzung}

Die Lizenz wird in \texttt{\%LOCALAPPDATA\%\textbackslash MaterialManager\_R03\textbackslash .license} gespeichert. Das Erststart-Datum wird mit einem SHA256-Hash gegen Manipulation geschützt.

% ==========================================
% 10. SYSTEMANFORDERUNGEN
% ==========================================
\section{Systemanforderungen}

\begin{table}[h]
\centering
\begin{tabularx}{\textwidth}{l X}
\toprule
\textbf{Anforderung} & \textbf{Minimum} \\
\midrule
Betriebssystem & Windows 10 (x64) oder höher \\
RAM & 4 GB \\
Festplatte & 200 MB (Self-Contained) \\
.NET & Nicht erforderlich (mitgeliefert) \\
Bildschirm & 1200$\times$720 px empfohlen \\
Netzwerk & Optional (für Multi-PC-Betrieb) \\
\bottomrule
\end{tabularx}
\caption{Systemanforderungen}
\end{table}

% ==========================================
% 11. FAZIT
% ==========================================
\section{Fazit}

\begin{successbox}[Zusammenfassung]
MaterialManager R03 ist eine \textbf{solide, praxistaugliche Branchenlösung} für die Blechlager-Verwaltung. Die Software bietet:

\begin{itemize}
    \item Vollständige Materialverwaltung mit automatischer Regalzuweisung
    \item Verschnitt-Optimierung und Reste-Management
    \item ERP-Kompatibilität (SAP, SAGE, BMEcat)
    \item Netzwerk-Betrieb für mehrere Arbeitsplätze
    \item Moderne Benutzeroberfläche mit Dark Theme
    \item Keine Installation nötig (USB-Stick-fähig)
\end{itemize}

Die Software ist bereit für den \textbf{produktiven Einsatz} in kleinen und mittleren Metallbau-Betrieben.
\end{successbox}

\vspace{1cm}

\begin{center}
\textcolor{gray}{\small Erstellt am \today\ $\bullet$ MaterialManager R03 v1.0.0}
\end{center}

\end{document}
